\mode*
\section{Method}

This study has two parts, mirroring the companion paper
\autocite{VTProgMisconceptions}.
The first is a \emph{literature review} that maps learners' misconceptions
and difficulties with the shell, the file system, composition, and remote
use, answering \cref{rq:misconceptions}:

\rqmisconceptions*

The second is an \emph{analysis} of those misconceptions through the lens of
variation theory: the misconceptions are treated as phenomenographic
observations from which we identify critical aspects
(\cref{rq:aspects}) and construct patterns of variation
(\cref{rq:patterns}):

\rqaspects*
\rqpatterns*

\Cref{rq:misconceptions} is answered by the review directly;
\cref{rq:aspects,rq:patterns} are answered analytically, and their answers
are validated by argument (inference to the best explanation), not yet
empirically.
The instruments of \cref{sec:instruments} prepare the empirical turn: they
both test the analytically derived aspects and collect the students' own
accounts, from which aspects can be found rather than argued.

\subsection{Literature review}
\label{sec:method-litreview}

The review has two phases.
The first phase is a seed round: a claim-driven search conducted for the
course-development work that this paper grows out of, with every retained
reference verified against at least its abstract and recorded with full
provenance (query, selection rationale, verbatim supporting quote,
verification level); the protocol and queries are in \cref{app:protocol}.
The second phase is a systematic, tool-supported multi-database search with
the \texttt{scholar} tool, as in the companion papers
\autocite{VTProgMisconceptions,VTDebug}; it covered shell and command-line
misconceptions, command-line education, learners' file-system mental models,
and SSH and remote-shell learning, and is documented in \cref{app:protocol}.

\subsection{Variation-theory analysis}
\label{sec:method-vt-analysis}

For each documented misconception we ask: which aspect had the learner not
discerned?
Aspects recurring across misconceptions are candidate critical aspects.
For each critical aspect we then construct patterns of variation---%
contrast, generalisation, fusion---that make the aspect discernible,
following the procedure of the companion paper
\autocite{VTProgMisconceptions}.
The shell makes this concrete: each pattern is realised as a sequence of
commands whose output the learner observes, varying one thing at a time
(contrast), varying the irrelevant surroundings (generalisation), or
varying several aspects at once (fusion)---the vary-and-observe loop that
programming environments afford for learning
\autocite{Svensson2020Affordances} and that variation theory has been shown
to support in lab-based computing teaching
\autocite{ThuneEckerdal2009Variation,ThuneEckerdal2018Lab}.

\subsection{Course instruments}
\label{sec:instruments}

The candidate critical aspects (\cref{rq:aspects}) are, at this stage, the
product of the analysis, not yet an empirical finding: we argue that a
learner who has not discerned an aspect will make the corresponding
mistake, but we have not measured whether discerning it is what the module
changes.
To make that future validation possible, and to let the students' own
ways of seeing correct the analysis, we prepare a pre/post-test that
collects two kinds of data.
It is a pair of diagnostic quizzes, run at the start and again at the end
of the terminal module, set up in the learning platform by the literate
program of \cref{app:quiz}.

The first kind is closed: one graded item per candidate critical
aspect, two for a few (\cref{tab:aspects}),
whose distractors express the documented misconceptions and the plausible
behaviour of a learner who has not discerned the aspect.
This is the usual shape of a diagnostic instrument
\autocite{KaltakciGurel2015Diagnostics} and of the pre- and post-tests
with which a learning study checks whether the aspects it taught were the
critical ones \autocite[Ch.~7]{NCOL}; a shell-specific instrument of this
kind does not exist---the nearest is a concept inventory for operating
systems internals \autocite{WebbTaylor2014OSCI}---so we write one.
Following the principle that a test of discernment must not point out
the aspect to be discerned \autocite[p.~91]{NCOL}, each item shows a
concrete command and its behaviour and leaves the interpretation to the
student.
Administered at the start, the closed items measure which aspects the
students already discern; administered again at the end, the paired
difference per item measures what the module made discernible, and the
distractors chosen tell which misconception stood in the way.
The closed items therefore operationalise the analysis of
\cref{sec:shell,sec:filesystem,sec:composition,sec:ssh}: they can
confirm or refute the aspects we argued for, but they cannot find an
aspect we did not think of, since every alternative was written by us.

The second kind is open, and is there to find them.
Six ungraded items ask the students for their own account of a shared
stimulus---what the terminal is and what happens when something is
typed at it; what happens to a line of plain English typed at the
prompt; how they would get a file from one folder to another; what goes
on in a pipeline; what has happened after logging in to a server and how
results get back to the laptop; and, once logged in, what a listing
shows---the open twin of a closed item, so that the account and the
choice of the same student can be compared.
These accounts are phenomenographic data in the strict sense: the
students' own descriptions of how they experience the phenomenon
\autocite{Phenomenography}.
The only earlier phenomenographic study of Unix interviewed six students
and found a hierarchical outcome space---Unix as a set of commands, as a
tool for solving certain problems, as a professional computing
environment---whose levels tracked the students' marks
\autocite{DoyleLister2006Unix}; its interview prompts asked, among other
things, what the student understood by \enquote{file system},
\enquote{pipe} and \enquote{script}, which are our objects exactly.
We analyse the accounts the same way: pooled per item and read for
qualitatively different ways of experiencing the phenomenon, which are
sorted into categories of description and ordered into an outcome space.
The critical aspects then fall out of the comparison: what a more
powerful way of seeing discerns that a less powerful one does not is,
by definition, an aspect critical for the learners in the lower category
\autocite[Ch.~2]{NCOL}.
Accounts from the start and the end of the module show, in addition,
the movement of individual students between categories.
The coding is inductive and done by hand; a language model drafts a
category and a discerned/undiscerned note per account, kept in separate
columns, to speed the reading of a large cohort (\cref{app:quiz-analysis}).

The order of inference matters, and so does its separation.
The open accounts of the first cohort are exploratory: an aspect that
emerges from them has not been tested by that cohort's closed items,
which were written from the literature-derived aspects before the
accounts were read.
Such an aspect earns a closed item in the next cohort's quiz, where it
is tested on students whose accounts played no part in finding it; we
never validate an aspect against the same cohort's closed answers from
which it was read off.
For the same reason the open items come before the closed ones in the
quiz, so that the distractors---which spell out the candidate
conceptions---cannot seed the accounts (\cref{app:open-items}).
A later cohort can also be given a two-tier form of the closed items,
asking the student to justify the chosen alternative, once the categories
to expect are known; for the first cohort that would only confirm what
we put into the distractors.

The module also runs, copied by hand---quizzes included---into the CS
programme's parallel course; nothing is deployed there, but the analysis
fetches the results from both courses and keeps the cohorts separable by
course (\cref{app:quiz-analysis}).

The end quiz answers its own items, so that the diagnostic also teaches:
under each closed item it shows why the keyed alternative is right and
what each distractor assumes, and under each open item the ways of seeing
we expect for that stimulus, ordered from the first impression to the one
the module aims at---a provisional outcome space, posited from the
candidate aspects since no account has yet been read
(\cref{app:open-items-feedback}).
It is shown only after submission and only at the end of the module, and
the coder never sees it; once the first cohort's accounts are coded, the
empirical outcome space takes its place.

The written accounts have the reach of a whole cohort but not the depth
of a conversation.
We therefore plan a small interview sub-study after the module:
semi-structured phenomenographic interviews with six to ten volunteers,
whose prompts are the six open stems followed up in conversation, with
the three prompts of \citeauthor{DoyleLister2006Unix} on the file
system, the pipe and the script added for comparability.
The transcripts are analysed together with the written accounts into one
outcome space; the interviews are where a category seen only dimly in the
accounts is confirmed or dissolved, and they are the counterpart of the
think-aloud sessions in the companion debugging study \autocite{VTDebug}.

The pre/post-test is one of two data sources the follow-up will have.
The other is the course's own assessment: the terminal and SSH labs are
graded by checking the state the student produced on the shared server,
and each check operationalises a candidate critical aspect---path
errors, for instance, are direct observations of undiscerned
file-system aspects (\cref{sec:filesystem}).
Each check failure is thus a misconception observation collected by the
assessment the course already runs, over repeated attempts (the grading
runs all year), giving attempt trajectories: which aspects block
students first, and which persist.
The two sources are kept methodologically separate to avoid
circularity: codes derived from the assessment traces are never
correlated with those same traces; the pre/post-test is the independent
measure of discernment.

\subsection{Ethics}
\label{sec:ethics}

Participation in the study is opt-in and requires informed consent; the
students' grades are unaffected by participation, and the data is
collected and analysed pseudonymously.
The consent item opens the pre-test (\cref{app:quiz}), in the same
wording as the course's other diagnostics, so a student who has answered
one recognises the next.
The consent covers the two quizzes---the closed items and the written
accounts alike---the course's own assessment record for the terminal and
SSH labs, which is the check results the grading produces anyway, over
repeated attempts, and being asked about the interviews.
The interviews themselves are asked for separately and recorded only
with explicit consent; the transcripts are pseudonymised at
transcription, and quotations are published only anonymised.
The parallel course runs the same module with the quizzes---consent item
included---copied over, so its students consent the same way, and the
analysis keeps the cohorts separable by course (\cref{app:quiz-analysis}).
Grading runs for all students regardless of consent; research extraction
happens afterwards and filters by consent (\cref{app:quiz-analysis}), so
no teacher sees during the course who consented.
The analysis identifies students by login ID rather than by name
(\cref{app:quiz-analysis}); the coding sheet, where a human codes the
accounts, is the only place the name and the pseudonymised data meet.
The accounts are pre-coded by a language model into columns of their
own before a human coder verifies them (\cref{app:quiz-analysis}); the
model sees only anonymised text---the analysis program strips the
identifiers and redacts incidental sensitive content before the call and
joins the suggested codes back itself---and no account is sent to a
model before that step is in place.
The collection was assessed against the Swedish Ethics Review Act
\autocite{Etikprovningslagen} with the institution's checklist: it
processes only ordinary personal data, the teaching is not manipulated
for research, and there is neither an intervention nor an obvious risk
of harm, so the act does not apply and no approval from the Ethics
Review Authority is required; the institution's research-ethics advisor
confirmed the assessment on 2026-09-01.
The personal data is processed under the GDPR with the institution as
controller: the students receive written information on the purpose,
the data, its handling and their rights, consent is opt-in and can be
withdrawn at any time, and the analysis uses pseudonymised data only.
The cohort of autumn 2026 answered the consent item before the written
information existed; it is delivered to them at the end of the course
and by email, and from the next cohort it accompanies the consent item.
The assessment accompanies this paper's source as
\texttt{ethics-application-datintro26.md}.
